\documentclass[11pt]{article}
\usepackage[T1]{fontenc}
\usepackage{textcomp}
\usepackage[margin=0.75in]{geometry}
\usepackage{amsmath,amssymb,amsthm}
\usepackage{array,booktabs}
\usepackage{hyperref}
\setlength{\parindent}{0pt}
\newtheorem{theorem}{Theorem}
\theoremstyle{definition}
\newtheorem{definition}{Definition}
\title{Flow Proof Artifact: Ideal-zero-product}
\author{Generated by \texttt{flow doc proof}}
\date{Flow Proof Book}
\begin{document}
\maketitle
Multiplying zero by a ring element stays in an ideal.\\[0.5em]
\textbf{Source.} dummit-foote — *Abstract Algebra*, §7.3\\[0.5em]
\bigskip
\noindent{\large \textbf{Derived fact 1.} \textit{If 0 in I then r * 0 in I} \hfill Ideal \cdot multiplication \cdot zero product stays in the ideal}\par
\smallskip\noindent\textit{Coordinate: Ideal \cdot multiplication \cdot zero product stays in the ideal}\par
\smallskip\noindent\textit{Source: dummit-foote}\par
\smallskip\noindent\textit{Built on: absorbs ring elements on the left, for multiplication on Ideal, ideals contain zero, for membership on Ideal}\par
\medskip
\noindent\textbf{Goal.} If 0 in I then r * 0 in I\par
\noindent$\displaystyle \forall I \in Ideal \forall r \in Ring\quad r \cdot 0 in I$\par
\medskip
\noindent
\begin{tabular}{@{} >{\raggedright\arraybackslash}p{0.06\textwidth} >{\raggedright\arraybackslash}p{0.47\textwidth} >{\raggedleft\arraybackslash}p{0.41\textwidth} @{}}
\textbf{\#} & \textbf{Proof} & \textbf{Mathematics} \\
\midrule
\textbf{1.} & We prove that zero product stays in the ideal for multiplication on Ideal. & \\[0.45em]
\textbf{2.} & We split into exhaustive cases — the claim must hold in each one. & \\[0.45em]
\textbf{3.} & Case 1 (see step 2): suppose 0 in I. & \\[0.45em]
\textbf{4.} & We invoke the derived fact governing multiplication on Ideal: absorbs ring elements on the left, for multiplication on Ideal (instantiated for I, r, 0). & \\[0.45em]
\textbf{5.} & We invoke the derived fact governing membership on Ideal: ideals contain zero, for membership on Ideal (instantiated for I). & \\[0.45em]
\textbf{6.} & From step 3, step 4, and step 5, this implies r times 0 in I. Together with the other cases (step 3), the goal is discharged. Hence proven. & $\displaystyle r \cdot 0 in I$ \\[0.45em]
\bottomrule
\end{tabular}
\label{thm:Ideal-multiplication-zero_product_stays_in_the_ideal}
\par\medskip\hrule
\end{document}
